\documentclass[11pt]{article}
\usepackage[T1]{fontenc}
\usepackage{lmodern}
\usepackage[margin=1in]{geometry}
\usepackage{amsmath,amssymb,mathtools}
\usepackage{booktabs}
\usepackage{enumitem}
\usepackage{xcolor}
\usepackage[colorlinks=true,linkcolor=blue!45!black,urlcolor=blue!45!black]{hyperref}

\newcommand{\bx}{\mathbf{x}}\newcommand{\by}{\mathbf{y}}
\newcommand{\dd}{\mathrm{d}}
\newcommand{\Mpl}{M_{\rm pl}}
\newcommand{\BoxK}{\Box_{\rm Kerr}}

\title{\textbf{The sourced Klein--Gordon equation behind the $2\to2$ rates}\\[3pt]
\large self-interactions, self-gravity, black-hole recoil, and the limits we compute}
\author{}
\date{16 November 2026 (Section~\ref{sec:pin} added 18 November)}

\begin{document}
\maketitle

\noindent
Every $2\to2$ rate in the level-evolution code comes from one equation,
\begin{equation}
(\BoxK-\mu^2)\,\phi_4 = S(x), \label{eq:master}
\end{equation}
solved with the Kerr Green's function. The self-interaction rates and the three levels of self-gravity
(Newtonian, Newtonian with black-hole recoil, fully relativistic) differ only in the source $S$.
This note states what $S$ is in general and how it is reduced in each limit.
Conventions: signature $(-,+,+,+)$, $G=1/\Mpl^2$, $\alpha=GM\mu$, and $\Box=\nabla^\mu\nabla_\mu$ of the Kerr metric of mass $M$ and spin $a$.

\section{The equation}

The field is a real scalar with potential $V=\mu^2f_a^2\,[1-\cos(\phi/f_a)]$, minimally coupled to gravity:
\begin{equation}
\Box_g\phi = V'(\phi)=\mu^2\phi-\frac{\mu^2}{6f_a^2}\phi^3+\dots,\qquad G_{\mu\nu}[g]=8\pi G\,T_{\mu\nu}[\phi]. \label{eq:full}
\end{equation}
The cloud is a superposition of quasi-bound Kerr modes,
$\phi=\sum_N\big[c_N\phi_N+{\rm c.c.}\big]$, with $\phi_N=e^{-i\omega_Nt+im_N\varphi}R_N(r)S_N(\theta)$ and
$(\BoxK-\mu^2)\phi_N=0$. The occupation number of level $N$ is proportional to $|c_N|^2$.

Both nonlinearities in \eqref{eq:full} generate, at order $c_1c_2c_3^*$, a field component with
\begin{equation}
\omega_4=\omega_1+\omega_2-\omega_3,\qquad m_4=m_1+m_2-m_3 .
\end{equation}
This component is not a level of the cloud. It is forced, and it obeys \eqref{eq:master} with $S\propto c_1c_2c_3^*$:
two quanta in levels 1 and 2 scatter, one ends in level 3, and the other (``4'') carries the remaining energy and angular momentum.
If $\omega_4<\mu$ it cannot propagate to infinity and is absorbed by the hole (``BH'' channels); if $\omega_4>\mu$ it escapes (``Inf'' channels).

\paragraph{Solution.} Equation \eqref{eq:master} separates. With $\Sigma=r^2+a^2\cos^2\theta$, $\Delta=r^2-2GMr+a^2$ and the spheroidal harmonics
$S_{l m_4}(\theta)$ of mode 4, the radial function of each multipole obeys
\begin{equation}
\frac{\dd}{\dd r}\!\Big(\Delta\frac{\dd R_{4l}}{\dd r}\Big)+\Big[\frac{K^2}{\Delta}-\lambda_l-\mu^2r^2\Big]R_{4l}=T_l(r),\qquad
T_l(r)=\int\!\dd\Omega\;\Sigma\,S(x)\,S^*_{lm_4}(\theta)\,e^{i\omega_4t-im_4\varphi},
\end{equation}
with $K=(r^2+a^2)\omega_4-am_4$. It is solved with the two homogeneous solutions $R_{\rm in}$ (ingoing at the horizon) and
$R_{\rm up}$ (decaying for $\omega_4<\mu$, outgoing for $\omega_4>\mu$) and their Wronskian $W$:
\begin{equation}
A^{\rm BH}_l=\frac1W\int_{r_+}^{\infty}\!\dd r\,R_{\rm up}(r)\,T_l(r),\qquad
A^{\infty}_l=\frac1W\int_{r_+}^{\infty}\!\dd r\,R_{\rm in}(r)\,T_l(r).
\end{equation}
These are the amplitudes of the ingoing wave at the horizon and of the outgoing wave at infinity. The rate is the corresponding flux,
$\Gamma\propto(\omega_4-m_4\Omega_H)\sum_l|A_l^{\rm BH}|^2$ or $\Gamma\propto k_4\sum_l|A_l^\infty|^2$ with $k_4=\sqrt{\omega_4^2-\mu^2}$.
This part, the left-hand side of \eqref{eq:master}, is treated identically for all sources: it is always the Kerr operator.

\section{The source in general}

Collecting the terms of order $c_1c_2c_3^*$ in \eqref{eq:full}, with $g=g_{\rm Kerr}+h$ and $h$ the metric perturbation produced by the cloud,
\begin{equation}
\boxed{\;S=\sigma_{12}\Big\{-\frac{\mu^2}{f_a^2}\,\phi_1\phi_2\phi_3^*
\;-\;\delta\Box\big[h^{(23)}\big]\phi_1\;-\;\delta\Box\big[h^{(13)}\big]\phi_2\Big\}\;} \label{eq:S}
\end{equation}
where $\sigma_{12}=1$ for different incoming levels and $\tfrac12$ for identical ones. The first term is the self-interaction.
In the other two, $\delta\Box[h]$ is the first-order change of the wave operator,
\begin{equation}
-\delta\Box[h]\phi=h^{\mu\nu}\nabla_\mu\nabla_\nu\phi+\big(\nabla_\mu h^{\mu\nu}-\tfrac12\nabla^\nu h\big)\nabla_\nu\phi , \label{eq:dbox}
\end{equation}
and $h^{(b3)}$ is the metric perturbation sourced by the \emph{transition} stress tensor of levels $b$ and 3,
\begin{equation}
\delta G_{\mu\nu}\big[h^{(b3)}\big]=8\pi G\,T^{(b3)}_{\mu\nu},\qquad
T^{(b3)}_{\mu\nu}=\nabla_\mu\phi_b\nabla_\nu\phi_3^*+\nabla_\nu\phi_b\nabla_\mu\phi_3^*
-g_{\mu\nu}\big(\nabla\phi_b\!\cdot\!\nabla\phi_3^*+\mu^2\phi_b\phi_3^*\big). \label{eq:lin}
\end{equation}
$T^{(b3)}$ oscillates at $\omega_g=\omega_b-\omega_3$ with azimuthal number $m_b-m_3$. Physically: the interference of levels $b$ and 3
is an oscillating mass distribution; its gravitational field acts on the quanta of level $a$ and scatters them into mode 4.
There are two terms because either incoming level can play the role of $b$.

Equations \eqref{eq:S}--\eqref{eq:lin} are exact to first order in the cloud mass. Nothing has been assumed about velocities, and the hole is
dynamical: the solution of \eqref{eq:lin} contains its tidal response and its recoil.
Since all three terms multiply $c_1c_2c_3^*$, the amplitudes add coherently,
\begin{equation}
A=A_{\rm SI}\Big(\frac{\Mpl}{f_a}\Big)^{\!2}+A_{\rm SG}\quad\Longrightarrow\quad
\Gamma=\Gamma_{\rm SI}\Big(\frac{\Mpl}{f_a}\Big)^{\!4}+\Gamma_{\times}\Big(\frac{\Mpl}{f_a}\Big)^{\!2}+\Gamma_{\rm SG},
\end{equation}
which is the form stored in the rate tables ($\Gamma_{\rm SI}$, $\Gamma_{\rm SG}$, $\Gamma_\times$ at $f_a=\Mpl$, and $\kappa=\Gamma_\times/2\sqrt{\Gamma_{\rm SI}\Gamma_{\rm SG}}$).

\section{Limits of the gravitational source}

\subsection{Fully relativistic (``exact'')}
Solve \eqref{eq:lin} on Kerr as it stands. We do this in Lorenz gauge, where the bracket in \eqref{eq:dbox} vanishes and
\begin{equation}
S_{\rm SG}=\sigma_{12}\big\{h^{(23)\,\mu\nu}\nabla_\mu\nabla_\nu\phi_1+h^{(13)\,\mu\nu}\nabla_\mu\nabla_\nu\phi_2\big\},
\end{equation}
as a boundary-value problem in $(r,\theta)$ at frequency $\omega_g$: all ten components of $h$, regular at the future horizon, outgoing at infinity.
This contains retardation, the gravitomagnetic and spatial components, frame dragging and the response of the hole.
For transitions with $|m_b-m_3|=1$ the hole is displaced by the dipole of the cloud. Obtaining that part of the solution needs a
separate procedure, in which the equations are solved in the frame of the hole; it is written out in Section~\ref{sec:pin}.

\subsection{Post-Newtonian expansion}
For a bound state $v^2\sim GM/r\sim\alpha^2/n^2$, so \eqref{eq:lin} can be expanded in $v$. The first correction (1PN) is of relative
order $\alpha^2$. We have computed it for four channels with identical incoming levels, as a test; it is not used in any table.
The ingredients are the following.

\paragraph{Non-relativistic field and its stress tensor.} Write $\phi=(2\mu)^{-1/2}\big(\psi\,e^{-i\mu t}+{\rm c.c.}\big)$, with
$|\dot\psi|\sim\mu\alpha^2|\psi|$ and $|\nabla\psi|\sim\mu\alpha|\psi|$. Dropping the terms that oscillate at $2\mu$,
\begin{equation}
\begin{aligned}
T_{00}&=\mu|\psi|^2+\tfrac{i}{2}\big(\bar\psi\dot\psi-\psi\dot{\bar\psi}\big)+\tfrac{1}{2\mu}|\nabla\psi|^2, &\qquad
T_{0i}&=-j_i,\quad \mathbf j=-\tfrac{i}{2}\big(\bar\psi\nabla\psi-\psi\nabla\bar\psi\big),\\
T_{kk}&=\tfrac1\mu|\nabla\psi|^2-\tfrac{3}{4\mu}\nabla^2|\psi|^2 ,
\end{aligned}
\end{equation}
so that $T_{00}\sim\rho$, $T_{0i}\sim\rho v$ and $T_{ij}\sim\rho v^2$.

\paragraph{Metric.} In harmonic gauge and in the near zone, with $\bar h_{\mu\nu}=h_{\mu\nu}-\tfrac12\eta_{\mu\nu}h$,
\begin{equation}
\bar h_{\mu\nu}(\bx)=4G\!\int\!\dd^3y\,\frac{T_{\mu\nu}(\by)}{|\bx-\by|},\qquad
h_{00}=\tfrac12\big(\bar h_{00}+\bar h_{kk}\big),\quad h_{0i}=\bar h_{0i},\quad
h_{ij}=\bar h_{ij}+\tfrac12\delta_{ij}\big(\bar h_{00}-\bar h_{kk}\big).
\end{equation}
The leading terms are Newtonian, $h_{00}=-2\Phi$ and $h_{ij}=-2\Phi\,\delta_{ij}$. At the next order there are four effects:
the gravitomagnetic potential $h_{0i}=-4G\!\int\!j_i/|\bx-\by|$, sourced by the mass current; the stresses $T_{kk}$ as a source of $h_{00}$;
the spatial curvature $h_{ij}$ acting on the motion of the field; and retardation, which for a source oscillating at
$k_0=\omega_b-\omega_3$ replaces $1/|\bx-\by|\to1/|\bx-\by|-\tfrac12k_0^2|\bx-\by|$.

\paragraph{Effective interaction.} Eliminating $h$ gives a source--source interaction,
\begin{equation}
L_{\rm eff}=G\!\int\!\!\int\frac{1}{|\bx-\by|}\Big[\tfrac12T_{00}T'_{00}+\tfrac12\big(T_{00}T'_{kk}+T_{kk}T'_{00}\big)-2\,T_{0i}T'_{0i}\Big]
+\text{retardation}+O(v^4),
\end{equation}
with unprimed quantities at $\bx$ and primed ones at $\by$. The three terms are Newton (including the $O(v^2)$ part of $T_{00}$),
stresses and spatial curvature, and gravitomagnetism. One further term is of the same order, because $GM/r\sim v^2$: the exchanged
graviton propagates in the field of the hole, which linearised gravity on flat space does not contain,
\begin{equation}
L_{G^2}=-\frac{G^2M\mu^2}{2}\!\int\!\!\int|\psi(\bx)|^2|\psi(\by)|^2
\Big[\frac{1}{r_xr_y}+\Big(\frac1{r_x}+\frac1{r_y}\Big)\frac{1}{|\bx-\by|}\Big].
\end{equation}

\paragraph{Amplitudes.} For $1+1\to3+4$, with the transition bilinears $\rho^{BA}=\bar\psi_B\psi_A$,
$\mathbf j^{BA}=-\tfrac{i}{2}(\bar\psi_B\nabla\psi_A-\psi_A\nabla\bar\psi_B)$ and
$t^{BA}=\nabla\bar\psi_B\!\cdot\!\nabla\psi_A-\tfrac34\nabla^2(\bar\psi_B\psi_A)$,
\begin{align}
A_{\rm N}&=-G\mu^2\!\int\!\!\int\frac{\rho^{41}\rho^{31\prime}}{|\bx-\by|}, \qquad\qquad
A_{h_{0i}}=+4G\!\int\!\!\int\frac{\mathbf j^{41}\!\cdot\mathbf j^{31\prime}}{|\bx-\by|},\nonumber\\
A_{h_{ij}}&=A_{h_{00}[T_{kk}]}=-\frac G2\!\int\!\!\int\frac{\rho^{41}t^{31\prime}+t^{41}\rho^{31\prime}}{|\bx-\by|}, \qquad
A_{\rm ret}=+\frac{G\mu^2k_0^2}{2}\!\int\!\!\int\rho^{41}\rho^{31\prime}\,|\bx-\by|,\nonumber\\
A_{T_{00}}&=-\frac G2\!\int\!\!\int\frac{\rho^{41}(\nabla\bar\psi_3\!\cdot\!\nabla\psi_1)'+(\nabla\bar\psi_4\!\cdot\!\nabla\psi_1)\,\rho^{31\prime}}{|\bx-\by|},\\
A_{G^2}&=+G^2M\mu^2\!\int\!\!\int\rho^{41}\rho^{31\prime}\Big[\frac1{r_xr_y}+\Big(\frac1{r_x}+\frac1{r_y}\Big)\frac1{|\bx-\by|}\Big].\nonumber
\end{align}
Mode 4 is defined only through the Green's function of Section~1, so each amplitude is rewritten as $\int\!\bar\psi_4\,S_X$ by moving the
derivatives off $\psi_4$ (integration by parts, and the continuity equation $\nabla\!\cdot\mathbf j^{31}=i\mu k_0\rho^{31}$). With
$\Phi=\int\!\rho^{31}/|\bx-\by|$ and $\mathbf A=\int\!\mathbf j^{31}/|\bx-\by|$, for example,
\begin{equation}
S_{\rm N}=-G\mu^2\,\psi_1\Phi,\qquad
S_{h_{0i}}=-4iG\,\mathbf A\!\cdot\!\nabla\psi_1+2G\mu k_0\,\psi_1\Phi .
\end{equation}
These sources go on the right-hand side of \eqref{eq:master} exactly like the Newtonian one, and
\begin{equation}
A_{\rm SG}=A_{\rm N}\Big(1+\sum_X\delta_X\Big),\qquad\delta_X=\frac{A_X}{A_{\rm N}}\propto\alpha^2,\qquad
\Gamma_{\rm SG}\to\Gamma_{\rm SG}\,\Big|1+\sum_X\delta_X\Big|^2 .
\end{equation}

\paragraph{Size.} Coefficients $\delta_X/\alpha^2$ at small $\alpha$ ($\tilde a=0.95$; the $h_{ij}$ column counts twice in the sum, once for
$h_{ij}$ and once for $h_{00}[T_{kk}]$):
\begin{center}
\begin{tabular}{@{}lcccccc@{}}
\toprule
channel & $h_{0i}$ & $h_{ij}$ & $T_{00}$ & ret. & $G^2$ & sum\\
\midrule
$211\times211\to322+{\rm BH}$ & $-1.14$ & $+0.25$ & $+0.13$ & $+0.10$ & $-0.46$ & $-0.87$ \\
$322\times322\to211+\infty$ & $-1.10$ & $+0.21$ & $+0.15$ & $+0.12$ & $-0.30$ & $-0.72$ \\
$322\times322\to544+{\rm BH}$ & $-0.32$ & $+0.09$ & $+0.05$ & $+0.03$ & $-0.19$ & $-0.24$ \\
$544\times544\to322+\infty$ & $-0.33$ & $+0.08$ & $+0.06$ & $+0.05$ & $-0.11$ & $-0.17$ \\
\bottomrule
\end{tabular}
\end{center}
Gravitomagnetism is the largest term and reduces the attraction (co-moving mass currents repel). Two caveats apply to these numbers.
The individual terms depend on the gauge; only their sum, with the wavefunctions written in matching (harmonic) coordinates, is
physical. And this calculation holds the hole fixed, so it does not contain the recoil of Section~3.4. The fully relativistic
calculation of Section~3.1 has neither limitation and replaces it.

\subsection{Newtonian}
Keep only the leading term. With $\partial_t\phi\simeq-i\mu\phi$, $T^{(b3)}_{00}\simeq2\mu^2\phi_b\phi_3^*\equiv\rho^{(b3)}$ and
$-\delta\Box[h]\phi\simeq h^{00}\partial_t^2\phi=2\mu^2\Phi\,\phi$, so
\begin{equation}
S_{\rm SG}^{\rm N}=2\mu^2\sigma_{12}\big\{\Phi^{(23)}\phi_1+\Phi^{(13)}\phi_2\big\},\qquad
\nabla^2\Phi^{(b3)}=4\pi G\rho^{(b3)}\;\Rightarrow\;
\Phi^{(b3)}(\bx)=-2G\mu^2\!\int\!\dd^3y\,\frac{\phi_b(\by)\,\phi_3^*(\by)}{|\bx-\by|}. \label{eq:newton}
\end{equation}
This is the $\mu\Phi\psi$ term of the Schr\"odinger--Poisson system. The Poisson equation is solved in flat space and instantaneously
($\omega_g\to0$), by multipoles: with $\phi_b\phi_3^*\to R_bR_3^*\sum_Lc_LY_{LM}$, $M=m_b-m_3$,
\begin{equation}
\Phi^{(b3)}\propto\sum_L\frac{4\pi}{2L+1}\,c_L\,Q_L(r)\,Y_{LM},\qquad
Q_L(r)=r^{-L-1}\!\!\int_0^r\!\dd r'\,r'^{L+2}R_bR_3^*+r^{L}\!\!\int_r^\infty\!\dd r'\,r'^{1-L}R_bR_3^* .
\end{equation}
Comparing with the first term of \eqref{eq:S}, in units of $GM$ for lengths,
\begin{equation}
\frac{A_{\rm SG}}{A_{\rm SI}}=8\,\alpha^2\Big(\frac{f_a}{\Mpl}\Big)^{\!2}\frac{J_{\rm SG}}{J_{\rm SI}},\qquad
J_{\rm SI}\leftrightarrow\phi_1\phi_2\phi_3^*,\quad
J_{\rm SG}\leftrightarrow\tfrac12\Big[\phi_1\!\int\!\frac{\phi_2\phi_3^*}{|\bx-\by|}+\phi_2\!\int\!\frac{\phi_1\phi_3^*}{|\bx-\by|}\Big],
\end{equation}
where $J$ denotes the amplitude obtained by pushing the indicated source through the Green's function of Section~1.

Only the right-hand side is Newtonian. The operator on the left of \eqref{eq:master} is still the Kerr one, and the levels
$\phi_{1,2,3}$ are Kerr bound states for $n\le8$ (hydrogenic above), exactly as for the self-interaction rates.

\subsection{Newtonian with black-hole recoil}
Equation \eqref{eq:newton} treats the hole as a fixed centre. It is not: a transition density with a dipole moment
($|m_b-m_3|=1$) pulls the hole back and forth with acceleration $\mathbf{a}_{\rm BH}=G\!\int\!\dd^3y\,\rho^{(b3)}(\by)\,\by/|\by|^3$.
The levels are defined in the hole's rest frame, which is therefore not inertial, and the field feels the fictitious potential
$\mathbf{a}_{\rm BH}\!\cdot\bx$:
\begin{equation}
\Phi^{(b3)}(\bx)=-G\!\int\!\dd^3y\,\rho^{(b3)}(\by)\Big[\frac{1}{|\bx-\by|}-\frac{\bx\cdot\by}{|\by|^3}\Big].
\end{equation}
This is the ``indirect term'' of celestial mechanics (equivalently the mass-polarisation term $\mathbf p_i\!\cdot\!\mathbf p_j/M$ of the
three-body problem). It is \emph{not} suppressed by the mass ratio: the force on the hole is $\propto M$, so its acceleration is
independent of $M$. Only the dipole changes,
\begin{equation}
Q_1(r)\;\to\;Q_1(r)-r\!\int_0^\infty\!\dd r'\,R_bR_3^*=r^{-2}\!\!\int_0^r\!\dd r'\,r'^3R_bR_3^*-r\!\int_0^r\!\dd r'\,R_bR_3^* ,
\end{equation}
i.e.\ the uniform field of the cloud outside $r$ drops out. For $211\times211\to322$ and $322\times322\to211$ the recoil reverses the sign
of the amplitude and reduces it to about $0.6$ of the fixed-hole value ($A_{\rm NR}/A_{\rm N}=-0.64$ and $-0.58$ at $\alpha=0.1$),
so the rate drops to $\sim0.35$--$0.4$ and the interference with the self-interaction becomes destructive.
The exact equation \eqref{eq:lin} contains the recoil as well; how it is extracted is the subject of the next subsection.

\subsection{Recoil in the relativistic calculation: pinning the hole}\label{sec:pin}

\paragraph{Why a special procedure.} Equation \eqref{eq:lin} does not single out a frame, and its solution for $|m_b-m_3|=1$ includes
the motion of the hole, a displacement $\mathbf d(t)\propto e^{-i\omega_gt}$. Two things make a direct solve unusable.
(i) A displaced hole is the same hole in shifted coordinates, so a static displacement is a pure-gauge $l=1$ perturbation
and an exact zero mode of the linearised operator. At frequency $\omega_g$, with the boundary conditions of Section~3.1, it is a
\emph{near}-zero mode with eigenvalue $\propto\omega_g^2$; for a spinning hole the tilt of the spin axis is a second one, with
eigenvalue $\propto\omega_g$. Since $GM\omega_g\simeq\tfrac12\alpha^3|n_3^{-2}-n_b^{-2}|$ is $10^{-3}$ to $10^{-6}$, the amplitude of these modes, which
\emph{is} the recoil, is fixed only by terms of relative size $(GM\omega_g)^2$, and discretisation errors are amplified by $1/(GM\omega_g)^2$, i.e.\ $10^{6}$ to $10^{11}$.
(ii) The levels $\phi_{1,2,3}$ and mode 4 are Kerr modes centred on the hole. The source \eqref{eq:S} must therefore be evaluated in
coordinates in which the hole stays at the origin, not in an asymptotically inertial frame in which it oscillates.

\paragraph{The frame.} We solve in coordinates attached to the hole: the horizon stays at $r=r_+$ with no $l=1$ displacement and no
tilt. This frame is accelerated, so far from the hole $h$ does not vanish but contains the field of a uniformly accelerated frame,
$h_{tt}=2\,\mathbf a_{\rm BH}\!\cdot\bx$, which in the Newtonian limit is the fictitious potential $\mathbf a_{\rm BH}\!\cdot\bx$ of Section~3.4.
The acceleration $\mathbf a_{\rm BH}$ is not put in by hand; it is an unknown of the calculation, fixed by the condition that nothing but
gravity acts on the hole. Write the Lorenz-gauge-fixed operator of \eqref{eq:lin} as $\mathcal A$, so that the equation is $\mathcal A\,h=8\pi G\,T^{(b3)}$.

\paragraph{Step 1: pin the hole and hold it with a strut.} Impose two conditions at the horizon,
\begin{equation}
p_d[h]\equiv\oint_{r_+}\!\!h^{rr}\,Y^*_{1M}=0,\qquad p_t[h]\equiv\big(\text{odd-parity }l=1\text{ part of }h_{vA}\big)_{r_+}=0,\qquad M=m_b-m_3=\pm1,
\end{equation}
the first removing the displacement of the horizon, the second its rotation ($v$ is the ingoing time, $A=\theta,\varphi$). A solution
of $\mathcal A h=8\pi GT$ with decaying boundary conditions is unique and will not satisfy them, so two unknowns are added: a force $F_d$ and a
torque $F_t$ acting on the hole, represented by fixed $l=1$ sources $z_d$ (in the $vv$ equation) and $z_t$ (in the odd $vA$ equation)
localised at the horizon. This gives the bordered linear system
\begin{equation}
\begin{pmatrix}\mathcal A & z_d\;\;z_t\\[2pt] p_d & 0\\ p_t & 0\end{pmatrix}
\begin{pmatrix}h_0\\ F^0_d\\ F^0_t\end{pmatrix}=\begin{pmatrix}8\pi G\,T^{(b3)}\\ 0\\ 0\end{pmatrix}, \label{eq:bordered}
\end{equation}
which is well conditioned: the pins remove the two near-zero modes. Its solution $h_0$ is the field of the cloud around a hole that is
\emph{held in place by an external agent} (a ``strut''), and $F^0$ is the force and torque the strut has to supply, i.e.\ minus those
exerted by the cloud. This is the relativistic version of the fixed-centre calculation of Section~3.3; it is not yet the physical solution.

\paragraph{Step 2: the field of an accelerated hole.} Next, the same system without matter, for a hole that the strut forces to
accelerate (or to rotate). Take the analytic fields of a uniformly accelerated and of a uniformly rotating frame,
\begin{equation}
U_d:\;\;h_{tt}=2\,r\sin\theta\,e^{iM\varphi}\quad(\text{i.e. }2\,\mathbf a\!\cdot\!\bx,\ \mathbf a\propto\hat{\mathbf x}\pm i\hat{\mathbf y}),\qquad
U_t:\;\;h_{ti}=(\boldsymbol\Omega\times\bx)_i,\quad\boldsymbol\Omega\propto\hat{\mathbf x}\pm i\hat{\mathbf y},
\end{equation}
multiply them by a smooth window $W(r)$, equal to 1 for $r<R_1$ (the whole cloud) and 0 for $r>R_2$, and solve for the remainder $q'_k$:
\begin{equation}
q_k=W\,U_k+q'_k,\qquad
\mathcal A\,q'_k+F^k_d z_d+F^k_t z_t=-\delta G\big[W U_k\big],\qquad p_{d,t}[q'_k]=-p_{d,t}[W U_k],\qquad k=d,t,
\end{equation}
with the same bordered matrix as in \eqref{eq:bordered} (one factorisation serves all three solves). Each $q_k$ satisfies the pin
conditions and the vacuum equations up to the strut, $\delta G[q_k]+F^kz=0$: it is the field of a hole dragged along by a strut force
$F^k$, described in its own accelerated frame.

\paragraph{Step 3: release the hole.} The physical hole is free. Superpose and require the total strut force and torque to vanish,
\begin{equation}
\boxed{\;h=h_0+a_d\,q_d+a_t\,q_t,\qquad F^0+a_d\,F^d+a_t\,F^t=0\;} \label{eq:release}
\end{equation}
a $2\times2$ linear system for $a_d$ and $a_t$. Then $\delta G[h]=8\pi G\,T^{(b3)}$ holds with no external force, the horizon is pinned,
and $h$ decays outside the window: this is the solution of \eqref{eq:lin} in the frame of the freely recoiling hole. The coefficient
$a_d$ is the acceleration of the hole, determined by force balance. For Schwarzschild only the displacement is treated this way.

\paragraph{Step 4: the source.} $S_{\rm SG}$ is \eqref{eq:dbox} evaluated with the full $h$ of \eqref{eq:release}, the numerical part and the
analytic part $a_dWU_d+a_tWU_t$ together (the latter is not in Lorenz gauge, so the bracket of \eqref{eq:dbox} is kept for it).
In the Newtonian limit $h_0$ gives the fixed-centre potential of \eqref{eq:newton} and $a_d\,U_d$ gives $\mathbf a_{\rm BH}\!\cdot\bx$,
so $A\to A_{\rm NR}$, the Newtonian amplitude with recoil. If the hole is \emph{not} released ($a_d=a_t=0$) one gets $A_{\rm N}$ instead.

\paragraph{Remarks.} The result does not depend on the choice of $U_k$, of the window or of the strut sources (any choice with a
nonzero strut force gives the same $h$ inside the window; checked, changes $\lesssim0.2\%$). The solution differs from the one in an
asymptotically inertial frame by the gauge transformation generated by the displacement, $\xi^\mu=d^\mu(t)$. Under a gauge
transformation the source changes by $(\BoxK-\mu^2)(\xi^\mu\partial_\mu\phi_a)$, i.e.\ $\phi_4\to\phi_4+\xi^\mu\partial_\mu\phi_a$, which leaves the
fluxes unchanged only if $\xi$ vanishes at the horizon and at infinity. A translation of the hole does not, which is why the frame has
to be specified, and why it has to be the one in which the levels are defined.
Channels without a dipole ($|m_b-m_3|\neq1$) need none of this: the hole does not move, and the direct solve of Section~3.1 is used.

\paragraph{Test.} Both statements of Step~4 were checked with the production code (18 November, $\tilde a=0.95$; amplitude ratios of the
leading multipole). In the first column the hole is held by the strut, in the second it is released:

\begin{center}
\begin{tabular}{@{}lcccc@{}}
\toprule
channel & $\alpha$ & held: $A/A_{\rm N}$ & released: $A/A_{\rm NR}$ & $A_{\rm NR}/A_{\rm N}$ \\
\midrule
$211\times211\to322+{\rm BH}$ & 0.05 & 0.9997 & 1.0041 & $-0.640$ \\
                              & 0.10 & 0.9989 & 1.0165 & $-0.639$ \\
$322\times322\to211+\infty$   & 0.05 & 0.9996 & 1.0063 & $-0.587$ \\
                              & 0.10 & 0.9985 & 1.0252 & $-0.580$ \\
$322\times544\to211+\infty$   & 0.05 & 1.0001 & 0.9769 & $+0.104$ \\
                              & 0.10 & 1.0005 & 0.9123 & $+0.107$ \\
$422\times844\to211+\infty$   & 0.075 & 1.0009 & 0.9722 & $-11.09$ \\
\bottomrule
\end{tabular}
\end{center}
Held, the relativistic amplitude equals the Newtonian one without recoil to $0.1\%$; released, it equals the Newtonian one with
recoil up to a correction that vanishes as $\alpha^2$ ($322\times544\to211$: $0.991$, $0.977$, $0.949$ at $\alpha=0.03$, $0.05$, $0.075$),
also where the recoil removes $90\%$ of the amplitude or exceeds the direct term tenfold. Over all 743 computed channel points with
a dipole exchange, in the 574 where the recoil changes the Newtonian rate by more than a factor of two the relativistic rate lies
within $30\%$ of the recoil-included value in $95.6\%$ and within $30\%$ of the fixed-hole value in none.

\section{What is used where}

\begin{center}
\begin{tabular}{@{}lll@{}}
\toprule
 & source $S$ & bound states $\phi_{1,2,3}$ \\
\midrule
Self-interaction tables & first term of \eqref{eq:S} & Kerr \\
Self-gravity tables (all runs so far) & Newtonian with recoil & Kerr for $n\le8$, hydrogenic above \\
Exact (overlay tables, 2{,}488 channel points) & \eqref{eq:lin}, Lorenz gauge & Kerr \\
1PN & test only, four channels & hydrogenic \\
\bottomrule
\end{tabular}
\end{center}
In all rows mode 4 is propagated with the Kerr Green's function, for the leading multipole $l_4$ (Inf channels) or $l_4=0,2$ (BH channels).
The hydrogenic approximation above $n=8$ is the main approximation on the wavefunction side. Against the paper's fully relativistic
self-interaction tables (17,065 common channels), the rates agree within $30\%$ for $93$--$94\%$ of the channels whose highest level has
$9\le n\le20$ ($98\%$ when all levels are relativistic in both), and differ by more than a factor of two for $2\%$; above $n=20$ the
agreement is worse ($82$--$86\%$ within $30\%$). The outliers are channels whose radial overlap integrals cancel strongly, so that
the $O(\alpha^2)$ difference between hydrogenic and Kerr wavefunctions is amplified.

\paragraph{When the Newtonian limit holds.} Two conditions: (i) $v^2\sim\alpha^2/n^2\ll1$; (ii) the cloud lies in the near zone of the
exchanged graviton, $\omega_gr\ll1$, with $\omega_g\sim\mu\alpha^2\,|n_3^{-2}-n_b^{-2}|/2$ and $r\sim n^2/(\mu\alpha)$. Both give corrections of
relative order $\alpha^2$. Measured ratios of the exact to the Newtonian-with-recoil amplitude at $\tilde a=0.95$:

\begin{center}
\begin{tabular}{@{}lccc@{}}
\toprule
channel & $\alpha=0.05$ & $\alpha=0.1$ & $\alpha=0.2$ \\
\midrule
$211\times211\to322+{\rm BH}$ (dipole) & 1.004 & 1.017 & 1.069 \\
$322\times322\to211+\infty$ (dipole) & 1.006 & 1.025 & 1.108 \\
$322\times322\to544+{\rm BH}$ & 1.0001 & 1.0004 & 1.002 \\
$844\times844\to544+\infty$ & 1.0006 & 1.0023 & 1.0093 \\
$10{,}4{,}4\times10{,}4{,}4\to322+\infty$ & 1.0004 & 1.0015 & -- \\
\bottomrule
\end{tabular}
\end{center}
The rate changes by the square. Corrections are at the per-cent level or below except for the dipole channels, where the recoil has
already removed most of the leading term and the $O(\alpha^2)$ remainder is relatively larger: $+3$ to $+5\%$ in the rate at
$\alpha=0.1$, $+14$ to $+23\%$ at $\alpha=0.2$.

\end{document}
